\scanpage{55}
% L01-03
\[
\left\{\begin{aligned}
\underline{\sigma}_0(\underline{s}_0)&\overset{\unsure{\mathrm{df}|\mathrm{dfn}|\mathrm{déf}}}{=}\underline{s}_1,
&\underline{\sigma}_0(\underline{a}_0)&=\underline{a}_0,
&\underline{\sigma}_0(\underline{f}_0)&=f_0\\
\underline{\sigma}_1(\underline{s}_0)&=\underline{s}_0,
&\underline{\sigma}_1(\underline{a}_0)&\overset{\mathrm{df}}{=}\underline{a}_1,
&\underline{\sigma}_{\unsure{2|1}}(\underline{a}_0)&=\underline{a}_0\\
\underline{\sigma}_2(\underline{s}_0)&=\underline{s}_0,
&\underline{\sigma}_2(\underline{a}_0)&=\underline{a}_0,
&\underline{\sigma}_2(\underline{f}_0)&\overset{\mathrm{df}}{=}\underline{f}_1
\end{aligned}\right.
\]

% L04-08
La donnée de $\varphi$ régulière équivaut à la donnée de $\varphi_G:\underline{\unsure{G|C}}\mathrel{\unsure{-|\longrightarrow}}\mathrm{Aut}(E)$ et $\varphi_{\unsure{\mathscr{R}|\mathcal{R}|R}}(\underline{r}_0)=(s_0,a_0,f_0)\in\mathrm{Drap}_{0,1,2}(E)$, satisfaisant certaines conditions. La donnée de $\varphi_G$ à son tour équivaut à la donnée des
% L09
\[
\sigma_0,\sigma_1,\sigma_2\in\mathrm{Aut}(E)
\]
% L10
satisfaisant certaines conditions, qu'on va expliciter.

% L11-13
Notons qu'indépendamment de la condition de régularité et du \deleted{\unsure{choix}} choix de $\underline{r}_0\in\underline{\mathscr{R}}$, la donnée de $\varphi$ implique
% L14
\[
\varphi_{\mathscr{R}}:\mathscr{R}\mathrel{\unsure{-|\longrightarrow}}\mathrm{Drap}_{012}(E)
\]
% L15-20
qui, dans le cas régulier, doit être une $G$-application, et qui permet de récupérer $\varphi_S,\varphi_A$ et $\varphi_F$ par passage au quotient. Dans le cas \deleted{\unsure{non}} régulier, comme $\mathscr{R}$ est un $G$-torseur, la donnée de $\varphi_{\mathscr{R}}$ équivaut à la donnée de $\varphi_{\mathscr{R}}(\underline{r}_0)\deleted{\unsure{\in|=}}r_0=(s_0,a_0,f_0)$.
% archive sheet number: 54
